\documentclass[11pt]{article}
\usepackage[T1]{fontenc}
\usepackage{lmodern}
\usepackage{amsmath,amssymb,amsthm,mathtools}
\usepackage[margin=1in]{geometry}
\usepackage{booktabs,array}
\usepackage[colorlinks=true,linkcolor=blue,citecolor=blue,urlcolor=blue]{hyperref}
\allowdisplaybreaks
\setlength{\emergencystretch}{2em}
\numberwithin{equation}{section}
\newtheorem{theorem}{Theorem}[section]
\newtheorem{proposition}[theorem]{Proposition}
\newtheorem{lemma}[theorem]{Lemma}
\newtheorem{corollary}[theorem]{Corollary}
\theoremstyle{definition}
\newtheorem{importedinput}[theorem]{Imported input}
\theoremstyle{remark}
\newtheorem{remark}[theorem]{Remark}
\DeclareMathOperator{\Rea}{Re}
\DeclareMathOperator{\Ima}{Im}
\DeclareMathOperator{\Log}{Log}
\DeclareMathOperator{\dist}{dist}
\newcommand{\dd}{\,d}
\newcommand{\HH}{\mathcal H}
\newcommand{\JJ}{\mathcal J}
\title{Microlocal Observations, Coherent Currents,\texorpdfstring{\\}{ }and Finite Zero Atlases for the Riemann Heat Flow}
\author{Drafted for Edward Baker}
\date{10 October 2026}
\begin{document}
\maketitle
\begin{abstract}
We organize phase space observations of the Riemann heat flow and their
finite arithmetic counterparts. A positive Husimi lift has an exact
generator with negative momentum diffusion, spatial drift, and a contact
term. Gaussian reconstruction retains these features in the observed
spatial jets; positivity alone does not exclude a real collision. A doubled
arithmetic density matrix retains both phase differences and reflected
phase sums. Centered logarithmic moments give exact block observations and
physical derivatives, while a complete coherent current yields a paid
necessary condition at a joint value--derivative zero. We also derive its
sharp support bound on the correlated first-jet error body. Under an
explicit imported holomorphic approximation estimate, outward interval
certificates prove joint nonvanishing and exactly 13, 12, 12, and 13 simple
heat zeros at every allowed time in four separate closed height windows.
Their proofs retain every natural-cutoff summand, physical spatial and
time transport, and the full value and derivative error payments. The
finite results do not cover the intervening heights or give a bound
uniform as time tends to zero. The remaining signed arithmetic estimate
is related to the density-strengthened threshold test and to the
threshold-preserving leading stationary reflection of the prime-torus reduction.
\end{abstract}
\tableofcontents

\section{Problem and scope}
The genuine scalar function is
\begin{equation}
 H_t(z)=\int_0^\infty e^{tu^2}\Phi(u)\cos(zu)\dd u,
 \quad
 \Phi(u)=\sum_{n\ge1}(2\pi^2n^4e^{9u}-3\pi n^2e^{5u})
                         e^{-\pi n^2e^{4u}}.
 \label{eq:heat}
\end{equation}
It satisfies $H_0(z)=\xi(1/2+iz/2)/8$ and
$\partial_tH_t=-\partial_z^2H_t$ \cite{polymath}.
The question motivating this investigation is whether its actual
arithmetic correlations supply an error-paid exclusion of
$H_t(x)=H_t'(x)=0$ at large real heights. Here a prime means a spatial
derivative at fixed time. It never means differentiation along a scaling
curve.

The continuous lift and the finite arithmetic approximant are different
objects. The lift represents \eqref{eq:heat} exactly. The finite sum
approximates a nonvanishing normalization of it through
Input~\ref{input:disk}. The exact finite algebra does not require that
input; all finite-certificate conclusions about genuine heat zeros do. The certificates
verify arithmetic signs and coverage after importing the analytic
approximation interface. They do not independently establish that
interface.

The principal finite conclusion is the following. Its quantitative data
and proof are given in Section~\ref{sec:certificates}.
\begin{theorem}[Four finite windows, under Input~\ref{input:disk}]
\label{thm:main}
For $M\in\{22066,22067,22068,22080\}$ put
\[
 t_M=\frac1{2\log M},\qquad x_M=4\pi M^2,\qquad
 \mathcal R_M=[t_M,1/20]\times[x_M,x_M+8].
\]
On each entire closed rectangle, the genuine pair $(H_t,H_t')$ has no
joint zero. For each allowed time, the numbers of zeros in its closed
height interval are respectively $13,12,12,13$, and every zero is simple.
With $Q_t=H_t/A_t$ and $L=\log(x/(4\pi))$, the corresponding uniform
lower bounds for $\|(Q_t,Q_t'/L)\|_2$ are
$0.03728,0.03117,0.02205,0.0032919$, each strict.
\end{theorem}
These are four disconnected windows, each with constant natural cutoff.
They are not complete cutoff cells. In their common time interval
$[t_{22066},1/20]$, their union contains exactly fifty simple zeros at
each time. This count concerns only that union. No assertion about zeros
in its gaps follows.

The analytical results explain why the finite certificates preserve the
information needed for this conclusion. Husimi positivity observes an
oscillatory Fourier transform. A positive rank-one arithmetic matrix can
have zero observed coordinates. A genuine macroscopic block can have
negative phase current. Its complement and interference can reverse that
sign. None of these facts decides the complete current on every genuine
candidate state. This paper consolidates the five microlocal notes
\cite{micro1,micro2,micro3,micro4,micro5}, with the signed-transform
background \cite{torus,reflection}; no literature novelty is asserted.

\section{Exact Wigner and Husimi observations}
\label{sec:lift}
Let $\Phi_e(u)=\Phi(|u|)$ denote the smooth even theta kernel in
\eqref{eq:heat}.
Its evenness follows from the theta transformation. For $u\ge0$, every
summand is positive since $2\pi n^2e^{4u}>3$; evenness gives positivity
on the whole real line. Define
\begin{equation}
 \psi_t(u)=e^{tu^2/2}\sqrt{\Phi_e(u)},\qquad
 \mathfrak m_t(u)=|\psi_t(u)|^2.
 \label{eq:wave}
\end{equation}
The theta tails are super-exponential. Their derivatives, and the
logarithmic derivatives of their positive square root, show that
$\psi_t\in\mathcal S(\mathbb R)$, with every finite time derivative
on compact real time intervals. In particular all polynomial
multiplication domains below contain the genuine trajectory.

With the convention
\begin{equation}
 W_t(u,p)=\frac1{2\pi}\int_{\mathbb R}e^{-ips}
     \psi_t(u+s/2)\overline{\psi_t(u-s/2)}\dd s,
 \label{eq:wigner}
\end{equation}
the Wigner function is real and Schwartz. Its momentum marginal is $\mathfrak m_t$.
Direct differentiation gives
\begin{equation}
 \partial_tW_t=u^2W_t-\frac14\partial_p^2W_t,
 \qquad H_t(z)=\frac12\int_{\mathbb R^2}e^{izu}W_t(u,p)\dd u\dd p.
 \label{eq:wigner-gen}
\end{equation}
Indeed the time multiplier inside \eqref{eq:wigner} is
$u^2+s^2/4$, whose Fourier image is $u^2-\partial_p^2/4$.
Integrating its momentum derivative gives zero boundary flux. The
negative scalar heat sign is therefore retained by this lift. Wigner
positivity is not assumed.

\begin{proposition}[Positive Husimi trajectory and its full generator]
\label{prop:husimi}
Fix $a>0$, $b_p=1/(16a)$, and set
\begin{equation}
 \HH_t=e^{a\partial_u^2+b_p\partial_p^2}W_t.
 \label{eq:husimi}
\end{equation}
This is a nonnegative function. On the genuine trajectory its exact
evolution is
\begin{equation}
 \partial_t\HH_t=\left(u^2+4au\partial_u+2a
                   +4a^2\partial_u^2-\frac14\partial_p^2\right)\HH_t.
 \label{eq:husimi-gen}
\end{equation}
\end{proposition}
\begin{proof}
For $g(y)=(\pi\sigma_g^2)^{-1/4}e^{-y^2/(2\sigma_g^2)}$,
$\sigma_g^2=4a$, the Wigner density is
$\pi^{-1}e^{-u^2/\sigma_g^2-\sigma_g^2p^2}$. It is exactly the
convolution kernel of \eqref{eq:husimi}: its variances are $2a,2b_p$.
Expanding the overlap integral gives
\[
 \HH_t(u,p)=\frac1{2\pi}\left|
 \int_{\mathbb R}\psi_t(y)
       \overline{e^{ip(y-u/2)}g(y-u)}\dd y\right|^2\ge0.
\]
Also $\int\HH_t=\|\psi_t\|_2^2$.
For Gaussian convolution,
$e^{a\partial_u^2}(uf)=(u+2a\partial_u)e^{a\partial_u^2}f$.
Applying this twice to \eqref{eq:wigner-gen} gives
$(u+2a\partial_u)^2=u^2+4au\partial_u+2a+4a^2\partial_u^2$.
Momentum convolution commutes with $\partial_p^2$. All operations can
be justified by convolution and integration by parts on Schwartz
functions, without a bounded inverse smoothing operator.
\end{proof}
The first-order term in \eqref{eq:husimi-gen} is spatial drift, and
$2a$ is its contact term. Negative momentum diffusion remains. The
proposition proves positivity of the actual smoothed trajectory; it
does not give positivity-preserving evolution on arbitrary nonnegative
phase space data. For general background on Husimi observations see
\cite{trushechkin}; the quantum evolution there is different from the
heat evolution here.

\subsection{Reconstruction and the spatial jet dictionary}
Define the smoothed scalar observation
\begin{equation}
 R_t^a(z)=\frac12\int_{\mathbb R^2}e^{izu}\HH_t(u,p)\dd u\dd p.
\end{equation}
Its marginal is the $u$-Gaussian convolution of $\mathfrak m_t$, so
\begin{equation}
 R_t^a(z)=e^{-az^2}H_t(z),\qquad
 H_t(z)=e^{az^2}R_t^a(z).
 \label{eq:reconstruction}
\end{equation}
The super-exponential marginal and its Gaussian convolution permit
entire extension of these identities. Thus the positive observation is
still oscillatory, and multiplication by a nonvanishing Gaussian
does not change joint value--derivative vanishing.

Put $p_k(x)=e^{-ax^2}\partial_x^ke^{ax^2}$. The first rows are
\begin{align}
 p_0&=1,&p_1&=2ax,&p_2&=2a+4a^2x^2,\\
 p_3&=12a^2x+8a^3x^3,&
 p_4&=12a^2+48a^3x^2+16a^4x^4.\nonumber
\end{align}
The recurrence $p_{k+1}=p_k'+2axp_k$ and the product rule give
\begin{equation}
 H_t^{(j)}=e^{ax^2}\sum_{k=0}^j\binom jk p_{j-k}(R_t^a)^{(k)}.
 \label{eq:gaussian-jets}
\end{equation}
At a collision this specializes to
\begin{align}
 H_t''&=e^{ax^2}(R_t^a)'',\nonumber\\
 H_t'''&=e^{ax^2}\bigl((R_t^a)'''+6ax(R_t^a)''\bigr),\\
 H_t^{(4)}&=e^{ax^2}\bigl((R_t^a)^{(4)}+8ax(R_t^a)'''
                  +(12a+24a^2x^2)(R_t^a)''\bigr).\nonumber
\end{align}
The reduced evolution is
\begin{equation}
 \partial_tR_t^a=-(\partial_x+2ax)^2R_t^a.
 \label{eq:reduced-husimi}
\end{equation}
If $a=a(t)$ moves, the generator additionally contains
$a'\partial_u^2+b_p'\partial_p^2$, and
\eqref{eq:reduced-husimi} additionally contains $-a'x^2R_t^a$.

For the normalizer $A_t$ defined below, write $\beta_A=A_t'/A_t$.
The correctly normalized observation and its derivative are
\begin{equation}
 Q_t=\frac{e^{ax^2}}{A_t}R_t^a,\qquad
 Q_t'=\frac{e^{ax^2}}{A_t}\bigl((R_t^a)'+(2ax-\beta_A)R_t^a\bigr).
 \label{eq:normalized-husimi}
\end{equation}
Equivalently the smoothed normalizer is $A_te^{-ax^2}$. Higher normalized
jets require differentiating the whole prefactor in
\eqref{eq:normalized-husimi}. An arbitrary numerical error $\varepsilon$
in $R_t^a(x)$ reconstructs to $e^{ax^2}\varepsilon$ in $H_t(x)$;
derivative errors use all coefficients of \eqref{eq:gaussian-jets}.
Exact smoothing of a known scalar error is different: its Gaussian
factor cancels upon reconstruction. No phase space truncation error
small enough for the large-height collision test is established here.

\section{The physical finite orbit and the imported analytic interface}
\label{sec:orbit}
For $s=(1-iz)/2$, with principal logarithms near the large positive
spatial axis, put
\begin{align}
 m_0(s)&=\Log s+\Log(s-1)-\frac s2\log\pi
       +\log\frac{\sqrt{2\pi}}{16}
       +\left(\frac s2-\frac12\right)\Log\frac s2-\frac s2,\nonumber\\
 \alpha(s)&=m_0'(s)=\frac1{2s}+\frac1{s-1}
                        +\frac12\Log\frac{s}{2\pi},\qquad
 m_t(s)=m_0(s)+\frac t4\alpha(s)^2.\label{eq:normalizer-data}
\end{align}
Define
\begin{align}
 P_{t,N}(s)&=\sum_{n\le N}
 e^{t\log^2n/4-(s+t\alpha(s)/2)\log n},\nonumber\\
 A_t(z)&=e^{(m_t(s)+m_t(1-s))/2},\qquad
 \theta_t(z)=\frac{m_t(s)-m_t(1-s)}{2i},\\
 F_{t,N}(z)&=e^{i\theta_t(z)}P_{t,N}(s)
               +e^{-i\theta_t(z)}P_{t,N}(1-s),\qquad
 Q_t(z)=H_t(z)/A_t(z).\nonumber
\end{align}
These are analytic at fixed $N$ on the neighborhoods in use. On the real
axis $A_t>0$, $F_{t,N}$ and $Q_t$ are real, and
\begin{equation}
 H_t=H_t'=0\quad\Longleftrightarrow\quad Q_t=Q_t'=0.
 \label{eq:joint-equivalence}
\end{equation}

At a real center define $L=\log(x/(4\pi))$, $\kappa=tL$, and
\begin{equation}
 0<t\le1/20,\qquad 1\le\kappa\le2,\qquad
 N=\left\lfloor\sqrt{x/(4\pi)+t/16}\right\rfloor.
 \label{eq:sector}
\end{equation}
All spatial derivatives use the same fixed time and integer cutoff.
The local scale $L$ specifies a Cauchy radius at that center; it is not
differentiated as a coordinate factor.

\begin{importedinput}[Complete fixed-cutoff disk estimate]
\label{input:disk}
For each center in \eqref{eq:sector}, the analytic fixed-cutoff
approximant above satisfies
\begin{equation}
 |Q_t(z)-F_{t,N}(z)|\le\eta(t,x)
 \quad (|z-x|\le1/L),\qquad
 \eta(t,x)=5e^{-tL^2/16-L/4}.
 \label{eq:disk}
\end{equation}
This is the repository's full interface from Heat Notes 8 and 13
\cite{heatinterface}, deduced there from the effective approximation
\cite{polymath}. It includes the symmetric normalizer, reflected branch,
and all changes of natural cutoff within the disk. The fixed integer in
$F_{t,N}$ is the center cutoff throughout that disk.
\end{importedinput}
Cauchy's formula then gives
\begin{equation}
 |Q_t^{(j)}(x)-F_{t,N}^{(j)}(x)|\le j!L^j\eta(t,x).
 \label{eq:cauchy}
\end{equation}
In particular the checker may use outward upper bounds for $\eta$ and
$L\eta$. The expression in \eqref{eq:disk} is the chosen majorant,
not a measured approximation error.

Write $\alpha=\alpha_r+i\alpha_i$, $\alpha'=U+iV$, where this prime
is the complex $s$ derivative. At the real center the exact formulas are
\begin{align}
 \alpha_r&=\frac L2+\frac14\log(1+x^{-2})-\frac1{1+x^2},&
 \alpha_i&=\frac{3x}{1+x^2}-\frac12\arctan x,\nonumber\\
 U&=\frac{7x^2-5}{(1+x^2)^2},&
 V&=\frac{x(x^2+5)}{(1+x^2)^2},\label{eq:real-data}\\
 \theta_t&=-\pi+\frac14\arctan x+\frac x4(1-L)
          -\frac x8\log(1+x^{-2})+\frac t2\alpha_r\alpha_i.\nonumber
\end{align}
Let
\begin{align}
 \sigma&=\frac12+\frac t2\alpha_r,&
 w_n&=n^{-\sigma}e^{t\log^2n/4},& T&=\frac{x-t\alpha_i}{2},\nonumber\\
 q_n&=w_ne^{i(\theta_t+T\log n)},&
 c&=\frac12(1+tU/2),& d&=tV/4,\label{eq:qdata}\\
 \Omega&=\frac12\Rea\{\alpha(1+t\alpha'/2)\},&
 \mu&=\Omega/c,&\epsilon&=d/c.\nonumber
\end{align}
Then, on the real axis,
\begin{equation}
 S=\sum_{n\le N}q_n,\qquad F_{t,N}=2\Rea S,\qquad
 \gamma_n:=q_n'/q_n=-d\log n-i(\Omega-c\log n).
 \label{eq:physical-derivative}
\end{equation}
The frequency $T$ is not the heat time. All $q_n$ use one actual common
frequency and carrier; no prime phase can be selected independently.
Off the real axis the approximant is the two analytic branches above,
not a real-part operation on a varying complex argument.

The finite approximation has a heat residual, as calculated in
\cite{torus}. For $b_n=e^{m_0(s)-s\log n+t(\alpha-\log n)^2/4}$,
$h_n=\alpha-\log n$, and $B_s=1+t\alpha'/2$,
\begin{equation}
 (\partial_t-\tfrac14\partial_s^2)b_n
 =\tfrac14\{h_n^2(1-B_s^2)-\alpha'B_s-(t/2)h_n\alpha''\}b_n.
 \label{eq:finite-residual}
\end{equation}
This follows from $(\log b_n)_s=h_nB_s$. Already at $t=0$ the
multiplier is $-\alpha'/4$. The finite sum is not an invariant finite
Newman heat manifold. A spatial error disk by itself does not bound
the time derivative of its error.

\section{Complete interference and centered block observations}
\label{sec:coherence}
For the doubled column $v=(q_1,\ldots,q_N,\overline q_1,\ldots,
\overline q_N)^T$ set $D=vv^*$. Let $\ell_0$ have all entries one
and let $\ell_1=(\gamma_n;\overline\gamma_n)_{n\le N}$. Then
\begin{equation}
 F=\ell_0v,\qquad F'=\ell_1v,\qquad
 \|\mathcal V_N\|_2^2=\tfrac14\ell_0D\ell_0^*
                         +\frac4{L^2}\ell_1D\ell_1^*,
 \quad\mathcal V_N=(F/2,2F'/L).
 \label{eq:density-observation}
\end{equation}
The upper-left matrix block contains $q_n\overline q_m$, with
phase differences. The upper-right contains $q_nq_m$, with reflected
phase sums. Positivity and rank one of $D$ do not ensure a nonzero
observed row. Dropping its phase-sum blocks changes the real
value--derivative observation.

For $B=\{\lfloor N/2\rfloor+1,\ldots,N\}$ and its complete
complement $C$, put $S_E=\sum_Eq_n$, $W_E=4iS_E'/L$,
and $\mathcal V_E=(\Rea S_E,\Ima W_E)$. Product-to-sum gives
\begin{align}
 \|\mathcal V_E\|_2^2
 &=\tfrac12\{|S_E|^2+|W_E|^2+\Rea(S_E^2-W_E^2)\},\nonumber\\
 \|\mathcal V_N\|_2^2
 &=\|\mathcal V_B\|_2^2+\|\mathcal V_C\|_2^2
     +E_{BC}^-+E_{BC}^+,\label{eq:recombination}\\
 E_{BC}^-&=\Rea(S_B\overline S_C+W_B\overline W_C),&
 E_{BC}^+&=\Rea(S_BS_C-W_BW_C).\nonumber
\end{align}
Neither cross term has an automatic sign. The actual adjacent cutoff
packet at $x=4\pi M^2$, $t=\kappa/(2\log M)$ has retained diagonal
energy asymptotic to $2\cos^2(\pi/8)w_M^2$ but coherent energy
$O(w_M^2/M^2)$ \cite{heatinterface,micro1}. Thus a positive fixed
fraction of that diagonal is not a universal packet lower bound.
This obstruction is not a collision of $H_t$.

\subsection{An exact relative-block moment transform}
Put $\ell_N=\log N$, $\delta_n=\log(N/n)$,
$a_N=\sigma-t\ell_N/2$, and
\begin{equation}
 Z_k=\sum_{n\in B}\delta_n^k
       e^{a_N\delta_n+t\delta_n^2/4-iT\delta_n},\qquad k\ge0.
 \label{eq:Z}
\end{equation}
Each $\delta_n\in[0,\log2)$. Completing the logarithmic square gives
the exact identity $S_B=q_NZ_0$. Define
$\lambda_N=\Omega-c\ell_N$, $\gamma_N=-d\ell_N-i\lambda_N$,
and $g=d-ic$. Physical differentiation at fixed $t,N$ gives
\begin{equation}
 a_N'=d,\qquad T'=c,\qquad
 S_B'=q_N(\gamma_NZ_0+gZ_1).
 \label{eq:block-first}
\end{equation}
The auxiliary identity $Z_1=i\partial_TZ_0$ freezes $a_N,t,N$;
it cannot replace \eqref{eq:block-first} in a physical derivative.
The signed block coordinates follow from
\[
 W_B=q_N\left[
 \left(\frac{4\lambda_N}{L}-\frac{4id\ell_N}{L}\right)Z_0
 +\left(\frac{4c}{L}+\frac{4id}{L}\right)Z_1\right].
\]

For higher jets set $\Gamma_n=\gamma_N+g\delta_n=\gamma_n$ and
use its physical coefficient derivatives with $\delta_n$ fixed.
The Bell multipliers are
\begin{align}
 P_0&=1,&P_1&=\Gamma_n,&P_2&=\Gamma_n^2+\Gamma_n',\nonumber\\
 P_3&=\Gamma_n^3+3\Gamma_n\Gamma_n'+\Gamma_n'',\label{eq:bell}\\
 P_4&=\Gamma_n^4+6\Gamma_n^2\Gamma_n'
        +3(\Gamma_n')^2+4\Gamma_n\Gamma_n''+\Gamma_n'''.\nonumber
\end{align}
The recurrence $P_{j+1}=\Gamma_nP_j+P_j'$ proves
$q_n^{(j)}=q_nP_j$. Writing $P_j=\sum_{k\le j}p_{jk}\delta_n^k$
therefore gives
\begin{equation}
 S_B^{(j)}=q_N\sum_{k\le j}p_{jk}Z_k,\qquad
 f_j=F_{t,N}^{(j)}=2\Rea(S_B^{(j)}+S_C^{(j)}).
 \label{eq:complete-jets}
\end{equation}
This includes derivatives of $U,V,c,\Omega$, amplitude drift and
normalizer motion. Its complete genuine payment is still
\eqref{eq:cauchy}. Removing the complement has no new analytic
approximation theorem attached to it.

\section{Paid complete-current and correlated candidate tests}
\label{sec:current}
Define the current without dividing by an amplitude:
\begin{equation}
 \JJ=-\Ima(S'\overline S),\qquad S=u+iv_0,\quad S'=u'+iv_1,
 \qquad\JJ=u'v_0-uv_1.
 \label{eq:current}
\end{equation}
The symbols $v_0,v_1$ here are scalar invisible quadratures, not the
doubled vector of \eqref{eq:density-observation}. The block formula is
\begin{equation}
 \frac{\JJ_B}{w_N^2}
 =\lambda_N|Z_0|^2+c\Rea(Z_1\overline Z_0)
                         -d\Ima(Z_1\overline Z_0),
 \label{eq:block-current}
\end{equation}
and the complete recombination is
\begin{equation}
 \JJ=\JJ_B+\JJ_C-\Ima(S_B'\overline S_C+S_C'\overline S_B).
 \label{eq:current-recombination}
\end{equation}
The common carrier phase value cancels from the current, which uses
phase differences. Its spatial motion remains through $\lambda_N$.
In fact $S\mapsto e^{i\chi(x)}S$ changes the current by
$-\chi'(x)|S|^2$; only a fixed common rotation leaves it invariant.
The actual real energy in \eqref{eq:recombination}
additionally retains reflected phase sums. The current is a
candidate-conditioned test, not an identification of that energy
with a positive complex norm.

\begin{proposition}[Rectangular current necessity]
\label{prop:current}
Under Input~\ref{input:disk}, a genuine joint zero at the center requires
\begin{equation}
 |u|\le\eta/2,\quad |u'|\le L\eta/2,\qquad
 |\JJ|\le\frac\eta2(L|v_0|+|v_1|).
 \label{eq:current-rect}
\end{equation}
Consequently a strict reverse of the last inequality for the complete
physical sum excludes a joint zero.
\end{proposition}
\begin{proof}
The first two bounds follow from $Q_t=Q_t'=0$, \eqref{eq:cauchy},
and $F=2u$, $F'=2u'$. Insert them in \eqref{eq:current}.
For fixed actual $v_0,v_1$, the right side is exactly the support of
this rectangle for the linear functional $u'v_0-uv_1$.
\end{proof}
There is no assumption $S\ne0$. The current test remains valid when
the complex coherent sum vanishes.

The error disk supplies a stronger necessary condition
\cite{schur,torus}:
\begin{equation}
 \left(\frac{2u}{\eta}\right)^2+\frac{2|u'|}{L\eta}\le1.
 \label{eq:schur}
\end{equation}
To derive it, apply Schwarz--Pick at zero to
$(Q_t-F_{t,N})(x+\zeta/L)/\eta$. Real symmetry makes its value
and derivative real there, and $|g'(0)|\le1-|g(0)|^2$ gives
\eqref{eq:schur} at a joint zero. Equality or a constant boundary
function follows by a limiting argument.

\begin{proposition}[Sharp correlated support for the current]
\label{prop:curved-current}
Set $A=L|v_0|$, $B_1=|v_1|$, and define
\begin{equation}
 h(A,B_1)=
 \begin{cases}
 A+B_1^2/(4A),&A>0,\ B_1\le2A,\\
 B_1,&A=0\ \text{or}\ B_1\ge2A.
 \end{cases}
 \label{eq:support}
\end{equation}
At a genuine joint zero under the same input,
$|\JJ|\le(\eta/2)h(A,B_1)$. This is the sharp support bound on
the first-jet body \eqref{eq:schur}, and is no larger than the
rectangular bound $(\eta/2)(A+B_1)$.
\end{proposition}
\begin{proof}
Let $s=2u/\eta$ and $r=2u'/(L\eta)$. Their necessary body is
$s^2+|r|\le1$. For fixed actual invisible quadratures the maximum of
$|Lv_0r-v_1s|$ is
\[
 \max_{0\le y\le1}\{A(1-y^2)+B_1y\}.
\]
Signs of $r,s$ can align the terms, and $y=|s|$.
For $A>0$ the stationary maximum has $y=B_1/(2A)$ when this is
at most one; otherwise the maximum is at $y=1$. For $A=0$ it is
$B_1$. This proves \eqref{eq:support}, including both degenerate cases.
\end{proof}
This proposition is a derived support refinement in the present
manuscript, also recorded in \cite{micro6}. It is sharp for the stated necessary body, not a
characterization of realizable arithmetic states. None of the retained
certificates below requires this sharper current test. A finite current
margin there uses \eqref{eq:current-rect}; the newer first-jet predicate
uses \eqref{eq:schur} directly.

\subsection{A genuine negative block and its complete recombination}
At $M=N=22066$, $t=t_M$, $x=x_M$, all $11033$ relative-block
terms are retained. The directed-power recomputation in the complete
rectangle checker certifies \cite{micro2,micro3}
\begin{equation}
 -181.169106<\JJ_B/w_M^2<-181.169105.
 \label{eq:negative-block}
\end{equation}
This stops a nonnegative-current claim for every genuine relative
block, even with coherent interference and physical drift retained.
It does not assert a complete-state collision. The same point
certificate gives the deliberately wider strict enclosures
\begin{center}
\begin{tabular}{lr}
\toprule
Quantity at $(t_M,x_M)$ & Enclosure\\
\midrule
$\JJ_B$ & $(-0.000673643,-0.000673642)$\\
$\JJ_C$ & $(12.32324,12.32325)$\\
Cross current & $(-0.215813,-0.215811)$\\
Complete $\JJ$ & $(12.10675,12.10677)$\\
$\|\mathcal V_M\|_2^2$ & $(3.13574,3.13575)$\\
\bottomrule
\end{tabular}
\end{center}
The two cross energies are approximately
$E_{BC}^-=-0.0437599560$ and $E_{BC}^+=0.00922600702$;
their signs are opposite. The complete-current tolerance is below
$0.061241$, leaving a paid gap above $12.0455$. These data show
why the actual complement must be recombined before judging a current.

\section{Coherent transport with complete physical payments}
\label{sec:transport}
The atlas machinery encloses a whole time--height rectangle, not just
samples of its center. For a selected $M$ put
$q_n^0=q_n(t_M,x_M)$, $\delta_n=\log(M/n)$,
$\omega_n=\delta_n/2$, and introduce the auxiliary sum
\begin{equation}
 G_M(h)=\sum_{n\le M}q_n^0e^{-i\omega_nh}.
 \label{eq:G}
\end{equation}
It is a coherent comparison function; the following payments relate
it to the physical sum $S_M(t,x_M+h)$.

\subsection{Spatial residual}
At fixed $t_M$ and cutoff $M$ define
\begin{equation}
 r_n=\gamma_n+i\omega_n
 =-\frac{t_MV\log n}{4}
 -i\{\Omega-c\log M+(c-1/2)\delta_n\}.
 \label{eq:r}
\end{equation}
This is centered before taking its absolute bound, so the nearly
cancelling carrier terms do not incur separate large payments. If
$|r_n|\le R_n$ throughout $0\le h\le H$, integration gives exactly
\begin{equation}
 q_n(t_M,x_M+h)=q_n^0e^{-i\omega_nh}
                 \exp\left\{\int_0^h r_n(x_M+y)\dd y\right\}.
\end{equation}
The elementary inequality $|e^z-1|\le e^{|z|}-1$ then proves
\begin{align}
 |S_M(t_M,x_M+h)-G_M(h)|&\le\varepsilon_{x,0}
   :=\sum_nw_n^0(e^{HR_n}-1),\nonumber\\
 |S_M'(t_M,x_M+h)-G_M'(h)|&\le\varepsilon_{x,1}
   :=\sum_nw_n^0\{\omega_n(e^{HR_n}-1)+R_ne^{HR_n}\}.
 \label{eq:spatial-payment}
\end{align}
Both are complete sums. They retain amplitude drift, carrier
curvature and the physical first derivative.

\subsection{Time and mixed derivatives at fixed physical height}
For $\ell=\log n$ direct differentiation gives
\begin{align}
 \chi_n:=q_{n,t}/q_n
 &=\ell^2/4-\alpha_r\ell/2
                      +i\alpha_i(\alpha_r-\ell)/2,\nonumber\\
 \gamma_{n,t}&=-V\ell/4
       -i\{U(\alpha_r-\ell)-\alpha_iV\}/4,\label{eq:time-multipliers}\\
 S_t&=\sum_nq_n\chi_n,\qquad
 S_{xt}=\sum_nq_n(\gamma_{n,t}+\gamma_n\chi_n).\nonumber
\end{align}
The imaginary part of $\chi_n$ includes the common time carrier
$\partial_t\theta_t=\alpha_r\alpha_i/2$. The mixed row includes
the time derivative of the spatial multiplier. Neither derivative
is taken along $tL=1$.

If $W_n,C_n,D_n,A_n^{\rm mix}$ bound respectively
$|q_n|,|\chi_n|,|\gamma_n|,|\gamma_{n,t}|$ on the whole rectangle,
set
\begin{equation}
 B_{t,0}=\sum_nW_nC_n,\qquad
 B_{t,1}=\sum_nW_n(A_n^{\rm mix}+D_nC_n).
 \label{eq:time-payment}
\end{equation}
Time displacement at most $\Delta t$ costs
$\Delta t B_{t,0}$ and $\Delta t B_{t,1}$. These are absolute
budgets only for the small physical motion; the signed main sum remains
coherent.

\subsection{Signed midpoint jets and their real Taylor remainder}
At midpoint $h_*$ compute every signed jet
\begin{equation}
 g_j=G_M^{(j)}(h_*)=\sum_{n\le M}(-i\omega_n)^jq_n^0e^{-i\omega_nh_*},
 \quad 0\le j\le K-1,
 \quad B_K=\sum_{n\le M}w_n^0\omega_n^K.
\end{equation}
For the degree-$K-1$ polynomial
$P(h)=\sum_{j<K}g_j(h-h_*)^j/j!$, the real Taylor integral
remainder gives, on $|h-h_*|\le d_*$,
\begin{equation}
 |G_M-P|\le B_Kd_*^K/K!,\qquad
 |G_M'-P'|\le B_Kd_*^{K-1}/(K-1)!.
 \label{eq:taylor}
\end{equation}
The derivative version expands $G_M'$ to degree $K-2$.
Along the real segment every phase exponential has modulus one,
so there is no additional exponential loss in these remainders.
Every low-index and complementary summand is present in $g_j,B_K$.

Combining \eqref{eq:spatial-payment}, \eqref{eq:time-payment},
and \eqref{eq:taylor} gives the componentwise enclosure radii
\begin{equation}
 E_0=B_Kd_*^K/K!+\varepsilon_{x,0}+\Delta tB_{t,0},\qquad
 E_1=B_Kd_*^{K-1}/(K-1)!+\varepsilon_{x,1}+\Delta tB_{t,1}.
 \label{eq:E}
\end{equation}
Each complex modulus bound safely encloses each real component.
To transfer the observations to $Q_t,Q_t'$, add respectively
$\eta$ and $L\eta$ to the real finite-sum errors $2E_0,2E_1$.

For the height width eight, $K=64$, $h_*=4$, $d_*=4$.
Each adaptive leaf uses an algebraic translation of the same polynomial
to its own midpoint. Interval Horner evaluation then uses a small
displacement. This changes the expression of $P$, not its coefficients'
meaning or the global Taylor remainder. It preserves the cancellations
already present in the signed jets.

\section{Finite certified rectangles and exact zero counts}
\label{sec:certificates}
\subsection{The first simple-zero rectangle}
The complete checker first encloses
\[
 [t_{22066},t_{22066}+8\cdot10^{-6}]
       \times[x_{22066}+0.3,x_{22066}+0.4].
\]
It uses the degree-five midpoint polynomial ($K=6$) about $h_*=0.35$
and radius $d_*=0.05$. Under Input~\ref{input:disk}, its fully
paid bounds are \cite{micro3}
\begin{align}
 -0.337&<Q_t(x_{22066}+0.3)<-0.246,\nonumber\\
 0.851&<Q_t(x_{22066}+0.4)<0.941,\\
 10.20&<Q_t'(x)<13.46.\nonumber
\end{align}
Its complete current lies between $5.07$ and $11.10$; the current
gap of \eqref{eq:current-rect} exceeds $4.92$. Opposite endpoint
signs and strict monotonicity give exactly one simple genuine zero
at each allowed time and joint nonvanishing everywhere in this smaller
rectangle. This rectangle lies inside the larger $M=22066$ atlas.

\subsection{Exact domains and arithmetic hulls}
For the four domains $\mathcal R_M$ of Theorem~\ref{thm:main},
the lower time is the exact expression $t_M$, not a rounded decimal.
For $h\ge0$ and $t\ge t_M$, monotonicity and
$t_M(2\log M)=1$ prove the exact sector lower edge $tL\ge1$.
The outward upper bounds in the records give $tL<1.000243<2$.
They also give throughout all four domains
\begin{equation}
 M^2+0.0031242<x/(4\pi)+t/16<M^2+0.639745<(M+1)^2.
 \label{eq:cutoff-check}
\end{equation}
Thus $N=M$ everywhere on each real rectangle. The complex disks may
still meet a cutoff change, already paid by Input~\ref{input:disk}.
Rounded endpoint hulls extend microscopically beyond exact $t_M,x_M$;
they are computational enclosures. The sector assertion concerns the
exact domain, with its lower edge proved symbolically.

The uniform payments satisfy $\eta<0.009642$ on all four windows.
The $M=22066$ atlas has $L\eta<0.192864$; each newer window has
$L\eta<0.192860$. The complete transport bounds for the three newer
windows include
\begin{center}
\begin{tabular}{lr}
\toprule
Budget & Outward upper bound\\
\midrule
Value Taylor remainder & $1.486\cdot10^{-6}$\\
Derivative Taylor remainder & $2.3775\cdot10^{-5}$\\
Spatial value payment & $2.598\cdot10^{-7}$\\
Spatial derivative payment & $3.306\cdot10^{-7}$\\
$B_{t,0}$ & $2202$\\
$B_{t,1}$ & $2194$\\
$E_0$ & $0.026720$\\
$E_1$ & $0.026647$\\
\bottomrule
\end{tabular}
\end{center}
The older atlas has the stronger total bounds
$E_0<0.019732$, $E_1<0.019684$. All signs below include these
physical and polynomial payments as well as the full analytic payments.

\subsection{Acceptance rules and complete closed coverage}
For a closed height leaf covering the whole allowed time interval,
let $a_0,a_1$ be lower distances from zero of the physical component
intervals for $u=\Rea S$ and $u'=\Rea S'$. Let
$\eta_{\rm up}$, $(L\eta)_{\rm up}$, $L_{\rm up}$ be outward
upper bounds. Define paid component gaps
\begin{equation}
 g_0=2a_0-\eta_{\rm up},\qquad
 g_1=2a_1-(L\eta)_{\rm up}.
 \label{eq:component-gaps}
\end{equation}
A leaf with $g_0>0$ has no genuine value zero. For $M=22066$,
every remaining accepted leaf has $g_1>0$ and a strict reverse of
\eqref{eq:current-rect} for its complete current enclosure.
The successful paid derivative condition alone already excludes a
joint zero; the current test records additional complete signed information.
For the three newer cutoffs every remaining leaf has $g_1>0$ and
the direct correlated predicate
\begin{equation}
 \underline{\mathcal P}
 =\left(\frac{2a_0}{\eta_{\rm up}}\right)^2
                 +\frac{2a_1}{(L\eta)_{\rm up}}>1.
 \label{eq:schur-lower}
\end{equation}
The calculation uses outward arithmetic to form a lower bound for
the left side of \eqref{eq:schur}. In these particular records,
$g_1>0$ already implies the second term of
\eqref{eq:schur-lower} exceeds one. The Schur check is active but
redundant for exclusion. No extra arithmetic exclusion from the curved
body alone is demonstrated.

The algorithm bisects any unresolved height leaf. Every retained leaf
is closed and covers the full exact time interval. Exact dyadic
adjacency, both height endpoints, and total width eight are verified;
every final leaf passes an acceptance rule. There are no omitted
unresolved strips. Quantitative joint floors follow from
\begin{equation}
 |Q_t|\ge g_0\quad\hbox{on value leaves},\qquad
 |Q_t'|/L\ge 2a_1/L_{\rm up}-\eta_{\rm up}>0
       \quad\hbox{on derivative leaves}.
 \label{eq:joint-floor}
\end{equation}
Taking the least of these component margins bounds
$\|(Q_t,Q_t'/L)\|_2$ on the whole rectangle. This normalized pair
has different scaling from $\mathcal V_N=(F/2,2F'/L)$.

\begin{center}
\small
\begin{tabular}{rrrrrr}
\toprule
$M$ & Leaves & Value leaves & Remaining & Zeros/time & Joint floor\\
\midrule
22066 & 46 & 29 & 17 & 13 & $>0.03728$\\
22067 & 29 & 16 & 13 & 12 & $>0.03117$\\
22068 & 28 & 14 & 14 & 12 & $>0.02205$\\
22080 & 32 & 15 & 17 & 13 & $>0.0032919$\\
\bottomrule
\end{tabular}
\end{center}
On the seventeen remaining $M=22066$ leaves the minimum paid
complete-current gap exceeds $0.11934$. For the newer cutoffs the
minimum first-jet lower bounds exceed respectively
$4.2338,3.2875,1.3415$. Only two remaining leaves per newer cutoff
also pass the coarse complete-current test. The other current
enclosures are inconclusive; a negative lower enclosure for a gap
does not certify a negative actual current. Even the older atlas
does not assert current positivity on every value-pruned leaf.

\subsection{Exact counts and smooth branches}
Merge adjacent non-value leaves into connected height bands. Each
band is checked again at both endpoints for all allowed times, with
the full value payment. Their signs are strictly opposite. The paid
$Q_t'$ has one strict nonzero sign on the whole band. Therefore
each band has exactly one zero per time. Every point outside these
bands belongs to a value leaf, so no other zero is possible.

For $M=22066$ the thirteen fixed offset bands and derivative signs are
\begin{center}
\begin{tabular}{ll@{\qquad}ll}
\toprule
Band & Sign & Band & Sign\\
\midrule
$[0.25,0.375]$ & $+$ & $[0.875,1]$ & $-$\\
$[1.5,1.625]$ & $+$ & $[2,2.25]$ & $-$\\
$[2.625,2.875]$ & $+$ & $[3.25,3.375]$ & $-$\\
$[3.75,3.875]$ & $+$ & $[4.375,4.5]$ & $-$\\
$[5,5.125]$ & $+$ & $[5.625,5.75]$ & $-$\\
$[6.125,6.5]$ & $+$ & $[7,7.125]$ & $-$\\
$[7.5,7.75]$ & $+$ & & \\
\bottomrule
\end{tabular}
\end{center}
Every such band has $|Q_t'|>3.08$. The certificates for the other
three cutoffs retain all their individual endpoint and derivative
intervals as well. At any normalized zero,
$H_t'=A_tQ_t'\ne0$. The implicit-function theorem, uniqueness,
and fixed disjoint trapping bands give smooth simple-zero branches
through the entire certified time interval, with local continuation
at its endpoints. No further continuation is asserted.

\begin{proof}[Proof of Theorem~\ref{thm:main}]
The complete interval calculations in the retained records
\cite{micro4,micro5} establish \eqref{eq:cutoff-check}, every
acceptance inequality, the closed coverage assertions, and the
endpoint and monotonicity signs just described. Their arithmetic
contract and reproducible sources are specified in
Section~\ref{sec:replay}. Input~\ref{input:disk} transfers each
component gap and endpoint sign to the genuine normalized function.
Equation~\eqref{eq:joint-floor} gives the displayed joint floors.
The band argument gives exactly $13,12,12,13$ zeros at every
allowed time, and \eqref{eq:joint-equivalence} transfers
nonvanishing and simplicity to $H_t$. Strict bounds apply at shared
leaf boundaries and at $t=1/20$ because all enclosures are closed.
\end{proof}

\section{Threshold jets and the stationary-reflection limitation}
\label{sec:threshold}
The finite current test does not require an all-real zero set. A
threshold fourth-jet test does, and it is a separate route.
For an ordinary double zero at an all-real time, the canonical-product
and zero-count inputs in \cite{torus,threshold} give, at sufficiently
large $x$,
\begin{equation}
 2q_3^2-3q_2q_4-\gamma_{\rm count}q_2^2\ge0,
 \quad q_j=Q_t^{(j)}(x),\quad
 \gamma_{\rm count}=18(\log A_t)''+\frac9{x^2}
                              +\frac{9\log x}{2D_0^2},
 \label{eq:threshold}
\end{equation}
where $D_0=8\pi(4C_{\rm count}+1)$ and $C_{\rm count}$ is the
symbolic constant in the imported count error. No numerical value of
that constant is used by the finite atlases. Opposite-root pairing
and additional counted real roots produce the last two terms.
The earlier microlocal notes printed the unstrengthened coefficient;
the complete density-strengthened version is used here.

Assemble the exact $f_j$ from \eqref{eq:complete-jets} and set
$\delta_j=j!L^j\eta$. A direct product perturbation payment is
\begin{align}
 \Delta={}&4|f_3|\delta_3+2\delta_3^2
    +3(|f_2|\delta_4+|f_4|\delta_2+\delta_2\delta_4)\nonumber\\
 &+|\gamma_{\rm count}|(2|f_2|\delta_2+\delta_2^2).
 \label{eq:threshold-payment}
\end{align}
Thus an opposite inequality sufficient to exclude this particular
ordinary threshold double zero is
\begin{equation}
 2f_3^2-3f_2f_4-\gamma_{\rm count}f_2^2+\Delta<0.
 \label{eq:open-raw}
\end{equation}
No finite-current certificate above establishes this inequality
uniformly. The selected normalizer coefficient must be retained
when recomputing its payment. Higher multiplicities can make a
fourth-jet test vacuous and require additional tests.

For comparison with the signed prime-torus reductions, define
\[
 M_j=\sum_{n\le N}q_n(\log n-\mu)^j=X_j+iY_j,
 \qquad \Gamma_{\rm count}=\gamma_{\rm count}/c^2.
\]
After paying the physical Bell residuals, the corresponding leading
centered quadratic is
\begin{equation}
 \mathcal K=2Y_3^2+3X_2X_4-\Gamma_{\rm count}X_2^2.
 \label{eq:centered-K}
\end{equation}
The raw leading jets are $-2c^2X_2,2c^3Y_3,2c^4X_4$.
They yield $4c^6\mathcal K$ in \eqref{eq:threshold}.
The measured payment in \cite{torus} adds their complete Bell
residuals to the genuine Cauchy errors. They cannot be dropped
merely because $x$ is large.

Heat Note 22 \cite{reflection} explains an exact near symmetry of
the prescribed leading stationary coefficient. For positive real $v$,
put $w(v)=v^{-\sigma}e^{t\log^2v/4}$, $P=T/(2\pi)$,
$\rho(v)=\log v-\mu$, and $q(v)=w(v)e^{i(\theta_t+T\log v)}$.
Define the leading coefficient of logarithmic degree $j$ at the
original-index saddle $v=P/k$ by
\begin{equation}
 B_j(k)=e^{i\theta_t}\frac{\sqrt P}{k}w(P/k)\rho(P/k)^j
          e^{i\{T\log(P/k)-T-\pi/4\}}.
\end{equation}
It can be written exactly as
\begin{equation}
 B_j(k)=e^{i\Theta}e^{\delta_\sigma(2\log k-\log P)}
          \overline{q(k)}\{-\rho(k)+\Delta_\mu\}^j,
 \label{eq:reflection}
\end{equation}
where
\[
 \delta_\sigma=\sigma-(1/2+(t/4)\log P),\quad
 \Delta_\mu=\log P-2\mu,\quad
 \Theta=2\theta_t+T\log P-T-\pi/4.
\]
Negative saddle curvature gives the phase $-\pi/4$
\cite{stationary}. The identity follows by expanding
$\log^2(P/k)$, the center and the carrier. For the prescribed data
on the closed subsector $1\le\kappa\le3/2$ as $t\downarrow0$,
\[
 \delta_\sigma=O(tx^{-2}+t^2x^{-1}),\qquad
 \Delta_\mu=O(x^{-2}),\qquad \Theta+2\pi=O(x^{-1}).
\]
On a fixed macroscopic dual range
$c_1\sqrt P\le k\le c_2\sqrt P$, $0<c_1<c_2$, this compares the coefficient to
$(-1)^j\overline{q(k)}\rho(k)^j$ with error $O_j(w(k)/x)$.
It is not an endpoint-uniform approximation theorem for the full
stationary integral.

The limiting algebraic map $M_j\mapsto(-1)^j\overline M_j$
fixes even cosine and odd sine moments, hence fixes
\eqref{eq:centered-K}. It changes the lower drift coordinate
$Y_1+\epsilon X_1$ to $Y_1-\epsilon X_1$; that discrepancy must
be paid. Moreover the saddle sends $[A,B]$ to $[P/B,P/A]$,
so $[N/2,N]$ maps approximately to $[N,2N]$, not the same
natural cutoff. Resonant sublattices acquire their own shifted
weights, centers and carriers. Thus the reflection alone supplies
neither a complete-cutoff identification nor an opposite threshold
sign. The coupled primitive-frontier representation in \cite{torus}
overlaps a full product-channel representation. Use one complete
representation or a disjoint common-factor split, with the appropriate
endpoint payments, so that product terms are not counted twice.

\section{Reproducibility and the remaining uniform estimate}
\label{sec:replay}
The numerical sources live under
\path{13_microlocal_phase_space/numerics/} relative to the
Newman-collisions investigation. The source hierarchy is
\begin{center}
\small
\begin{tabular}{ll}
\toprule
Source & Role\\
\midrule
\path{check_phase_space.py} & Rational algebra and floating diagnostics\\
\path{check_block_current.py} & Initial genuine block interval source\\
\path{check_complete_current_rectangle.py} & Complete current and small rectangle\\
\path{check_candidate_current_atlas.py} & $M=22066$ closed atlas\\
\path{check_multi_cutoff_candidate_family.py} & Three newer cutoff atlases\\
\bottomrule
\end{tabular}
\end{center}
The first source has exact rational controls and floating diagnostic
checks; those diagnostics are not genuine arithmetic sign certificates.
The other sources use 60-digit outward Decimal intervals. Pi comes
from the alternating Machin arctangent identity with first-omitted-term
bounds. Large arctangent uses its reciprocal argument. Reduced
trigonometric polynomials have explicit Taylor remainders on their
checked argument range. Monotone logarithm and exponential endpoints
are enclosed by adjacent values around correctly rounded outputs
\cite{decimal}. No binary float enters a certificate assertion.

The original block source used \texttt{Context.power}. The later
complete-rectangle source replaces nonnegative integer interval
powers locally by repeated directed multiplication with exact
sign/parity and zero-crossing rules. Both atlas sources inherit this
override. The documentation gives a weaker rounding contract for
the C implementation of general power \cite{decimal}; therefore
the proof record for \eqref{eq:negative-block} here uses the later
complete recomputation. The preserved older files are not rewritten.
Rational endpoint/parity and polynomial-shift controls test the
relevant arithmetic routines.

Each later source verifies its predecessor's retained SHA-256 before
importing it. The build records identify those dependencies and the
small retained certificates. Hashes establish file identity; fresh
outward arithmetic decides the signs and coverage. The analytic
interface and the Decimal operation contracts remain explicit
dependencies of the computer-assisted theorem, not formalized proofs
of those contracts. Internal replays do not constitute independent
mathematical review.

From the repository root, reproducible fresh atlas outputs can be
requested without replacing retained records:
\begin{quote}\small
\path{python3 papers/quasi-rh-exponent-descent/newman_collisions/13_microlocal_phase_space/numerics/check_candidate_current_atlas.py}
\path{--output /tmp/microlocal-atlas-replay.json}

\path{python3 papers/quasi-rh-exponent-descent/newman_collisions/13_microlocal_phase_space/numerics/check_multi_cutoff_candidate_family.py}
\path{--M 22067 --output /tmp/microlocal-M22067-replay.json}
\end{quote}
Use the latter with $M=22068$ and $22080$ for the other windows.
The complete-rectangle source prints its JSON to standard output,
which can be redirected to a temporary file. All retained source and
certificate files satisfy the repository's small-file policy; no
large derived archive or manuscript snapshot folder is needed.

The additional \path{check_manuscript_certificate_summary.py} inspects
the retained records using exact rational arithmetic. It verifies twelve
build-record file identities, all 135 adjacent atlas leaves, their
conservative joint floors and Schur predicates, and fifty signed-monotone
zero bands. It does not replace the fresh interval-sum replays or prove
the imported analytic interface. Its small summary is retained as
\path{numerics/MANUSCRIPT_CERTIFICATE_SUMMARY_20261010.json}.

The finite successes identify a precise next obligation. Under the
same full-sum candidate constraints, one needs a complete-current
margin exceeding its invisible-quadrature support payment, or the
opposite paid threshold sign \eqref{eq:open-raw}, uniformly on a
specified growing family. The present time windows stay near
$t=1/20$. They neither approach zero time through a controlled
infinite family nor bridge the height gaps or cutoff boundaries.
Higher multiplicities, small-curvature cases, complementary
parameters and the time-zero endpoint remain separate obligations
for the global research program. Positivity of the continuous lift,
block signs, and the invariant leading stationary reflection do not
close them.

\section*{Acknowledgements and preparation record}
This manuscript was drafted for Edward Baker with substantial assistance
from OpenAI's Codex large language model. The model is identified in this
session as GPT-6 (Codex); the exact serving variant and configured
reasoning effort are not exposed and are not inferred. LLM assistance
included source synthesis, algebraic derivation, exposition and internal
cross-audits. Those checks are not independent mathematical validation.
Earlier source notes and certificates retain their own preparation
metadata. The genuine analytic approximation and classical heat
normalization are credited below. The sharp correlated current-support
formula is derived in Proposition~\ref{prop:curved-current}; the finite
certificates use their preserved acceptance rules.

\begin{thebibliography}{99}
\bibitem{polymath}
D. H. J. Polymath,
\emph{Effective approximation of heat flow evolution of the Riemann
$\xi$ function, and a new upper bound for the de Bruijn--Newman constant},
arXiv:1904.12438v2 (2019), Theorem 1.3.
\url{https://arxiv.org/abs/1904.12438}.

\bibitem{trushechkin}
A. S. Trushechkin,
\emph{Semiclassical evolution of quantum wave packets on the torus beyond
the Ehrenfest time in terms of Husimi distributions},
arXiv:1607.07572 (2016).
\url{https://arxiv.org/abs/1607.07572}.

\bibitem{decimal}
Python documentation, \emph{decimal---Decimal fixed-point and
floating-point arithmetic}, logarithm, exponential and power contracts.
\url{https://docs.python.org/3/library/decimal.html}.

\bibitem{stationary}
NIST Digital Library of Mathematical Functions, \S2.3(iv),
\emph{Method of stationary phase}.
\url{https://dlmf.nist.gov/2.3#iv}.

\bibitem{micro1}
Edward Baker research project, with LLM assistance,
\emph{Husimi readout drift and coherent interference}, microlocal Note 1,
10 October 2026; corresponding internal Review 1 and scout record.
\path{notes/1_HUSIMI_READOUT_DRIFT_AND_COHERENT_INTERFERENCE_20261010.md}.

\bibitem{micro2}
Same project, \emph{Centered relative block and an adverse genuine
coherent current}, microlocal Note 2, 10 October 2026; internal Review 2.
\path{notes/2_CENTERED_RELATIVE_BLOCK_AND_ADVERSE_COHERENT_CURRENT_20261010.md}.

\bibitem{micro3}
Same project, \emph{Complete current and a paid rectangle of genuine
simple heat zeros}, microlocal Note 3, 10 October 2026; internal Review 3.
\path{notes/3_COMPLETE_CURRENT_AND_PAID_SIMPLE_ZERO_RECTANGLE_20261010.md}.
Certificate: \path{numerics/COMPLETE_CURRENT_RECTANGLE_CERTIFICATE_20261010.json}.

\bibitem{micro4}
Same project, \emph{A complete candidate-current atlas and thirteen
simple heat-zero branches}, microlocal Note 4, 10 October 2026;
internal Review 4.
\path{notes/4_CANDIDATE_CURRENT_ATLAS_AND_THIRTEEN_SIMPLE_ZERO_BRANCHES_20261010.md}.
Certificate: \path{numerics/CANDIDATE_CURRENT_ATLAS_CERTIFICATE_20261010.json}.

\bibitem{micro5}
Same project, \emph{A finite multi-cutoff family with active correlated
candidate tests}, microlocal Note 5, 10 October 2026; internal Review 5.
\path{notes/5_MULTI_CUTOFF_CORRELATED_CANDIDATE_FAMILY_20261010.md}.
Certificates: \path{numerics/MULTI_CUTOFF_M22067_CERTIFICATE_20261010.json},
\path{numerics/MULTI_CUTOFF_M22068_CERTIFICATE_20261010.json},
\path{numerics/MULTI_CUTOFF_M22080_CERTIFICATE_20261010.json}.

\bibitem{micro6}
Same project, \emph{Correlated current support and manuscript synthesis},
microlocal Note 6, 10 October 2026.
\path{notes/6_CORRELATED_CURRENT_SUPPORT_AND_MANUSCRIPT_SYNTHESIS_20261010.md}.
The manuscript audit and fresh replay record are in
\path{reviews/6_MANUSCRIPT_INTERNAL_REVIEW_20261010.md}.

\bibitem{heatinterface}
Same project, \emph{Signed shrinking-time collision vector and phase
obstructions}, Heat Note 8, 9 October 2026, and
\emph{Recent heat results and paths forward}, Heat Note 13,
9 October 2026. Source directory: \path{../notes/}.
The fixed-cutoff disk deduction is in Note 8, Proposition 1,
and the parent \path{../newman_collision_reductions.tex}, Section 5.

\bibitem{schur}
Same project, \emph{Correlated holomorphic payments and candidate
coverage}, Heat Note 18, 10 October 2026.
\path{../notes/18_CORRELATED_HOLOMORPHIC_PAYMENTS_AND_CANDIDATE_COVERAGE_20261010.md}.

\bibitem{threshold}
Same project, \emph{Analytical Schur cones and density-strengthened
thresholds}, Heat Note 20, and \emph{Deflated local growth,
density-strengthened thresholds, and multiplicity}, Heat Note 21,
10 October 2026. Source directory:
\path{../notes/}.

\bibitem{torus}
Same project, \emph{Signed Arithmetic Reductions on the Prime Torus
for Riemann Heat-Flow Collisions}, manuscript, 10 October 2026.
\path{../09_prime_phase_torus/prime_phase_torus_signed_reductions.tex}.
Its Notes 1--9 and internal manuscript Review 10 supply the detailed
signed transformations, payments and limitations.

\bibitem{reflection}
Same project, \emph{Signed pair transforms and the prescribed stationary
reflection}, Heat Note 22, 10 October 2026.
\path{../notes/22_SIGNED_PAIR_TRANSFORMS_AND_STATIONARY_REFLECTION_20261010.md}.
\end{thebibliography}
\end{document}
